\scanpage{48}
qu'on peut interpréter comme l'équation de
Hamilton-Cayley pour la \inserted{l'op. de} translation \deleted{\unsure{induite}} définie
par $u$ dans $B$.

\uline{8} \uline{Cas où 2 inversible sur} $S$ (car résiduelle $\ne2$)

Dans ce cas la donnée d'une section $\underline b=\mathrm{tr}$ de
$2F$ i.e. d'une forme $\mathrm{tr}$ sur $B$ telle que $\mathrm{tr}(e_B)=2$,
équivaut, \deleted{\unsure{(en en divisant par} \inserted{en prenant} (en posant $\tfrac12\underline b=\tfrac12\mathrm{tr}\,u\unsure{:|=}\overset{\unsure{s.}}{\unsure{s}}$}
\[
 s_0=\tfrac12\deleted{\unsure{\mathrm{tr}|\underline b}}\underline b=\tfrac12\mathrm{tr}\,\unsure{u|\ldots} \tag*{(88)}
\]
à la donnée d'une section \inserted{origine} $s_0$ de $F$, ou encore
d'un splittage de l'extension $B$ de $\check L$ par $\mathcal O_S$ en
un iso
\[
\alpha_{s_0}:\underline O_S+\check L\xrightarrow{\sim}B \tag*{(89)}
\]
Le sous-module $\alpha_{s_0}(\check L)$ de $B$, supplémentaire
de $\mathcal O_S$, n'est autre que le \uline{noyau} de la
forme $\mathrm{tr}$ (qui est aussi celui de la forme
$\tfrac12\mathrm{tr}$ \ldots), \inserted{ou encore le sous-espace propre de $\sigma$ pour val. propre $-1$.}
% slanted
\margin{NB en termes de $q=\unsure{\mathcal N|\mathscr N|N}_{B/\mathcal O_S}$, $\alpha_{s_0}(\check L)$ est l'\uline{orthogonal} de $e_B$, i.e. de $\mathcal O_S e_B$.}
Ceci posé, la donnée d'une
structure d'Algèbre sur $B$ compatible avec
$(e_B,\mathrm{tr})$ équivaut à la donnée de
$\underline b(s_0)$\inserted{$\in\Gamma L$}, $\underline c(s_0)$\inserted{$\in\Gamma L^{\otimes2}$} correspondant au scindage $s_0$.
Or \inserted{\unsure{(14))|(14)}} \deleted{\unsure{au}} $\underline b_{\unsure{0}}=\underline b(s_0)$ est la forme trace restreint à
$\alpha_{s_0}(\check L)$, et dire que \deleted{$\underline b(s_0)$} $s_0$ est le scindage canonique
signifie exactement que
\[
\deleted{\unsure{b=}}\underline b(s_0)=0,\tag*{(90)}
\]
% archive sheet number: 47
